% !PRINTED 576
\weberpage{576}
%<<CONT
according to the meaning (12) of $\Pi(t)$, it follows that:
\[ \Pi'(\tau) \equiv (\tau-\tau_1)(\tau-\tau_2) \cdots (\tau-\tau_{f-1}) \equiv 0 \pmod{\mathfrak p}. \]
But this is impossible, because the $\tau, \tau_1, \ldots \tau_{f-1}$, as we have seen, are incongruent with respect to the modulus $\mathfrak p$. With this, then, our Theorem 4. is completely proved. We can accordingly, if $x, y$ denote two new variables, form a functional $\pi$ of the ideal $\mathfrak p$:
\weq{23}{\pi = x p + y P(\tau).}
We now consider integral forms $\Phi(t)$ in $R$ which by the substitution $t=\tau$ become divisible not only by $\mathfrak p$, but by the natural prime $p$, and which therefore satisfy the congruence
\weq{24}{\Phi(\tau) \equiv 0 \pmod p}
Let the prime $p$ be decomposed into its prime factors in $\Omega$ in the following manner:
\weq{25}{p = \mathfrak p \mathfrak p_1 \mathfrak p_2 \cdots,}
where
\[ \mathfrak p, \mathfrak p_1, \mathfrak p_2, \ldots \]
are equal or different prime ideals of the degrees
\[ f, f_1, f_2, \ldots \]
To these prime factors there corresponds (by Theorem 3.) a sequence of integral rational forms
\[ P(t), P_1(t), P_2(t), \ldots \]
of the degrees $f, f_1, f_2, \ldots$, and if, say, $\mathfrak p$ is identical with $\mathfrak p_1$, then $P(t)$ is also identical with $P_1(t)$. If one takes the norm on the right and on the left in (25), and takes account of the formula $N(\mathfrak p) = p^f$ [\S.\,140, (3), \S.\,151, (3)], it follows that
\weq{26}{n = f + f_1 + f_2 + \cdots}
Now if $\Phi(t)$ is an integral rational form satisfying the condition (24), it follows from Theorem 3. that
\weq{27}{\Phi(t) \equiv P(t) \Phi_1(t) \pmod p,}
where $\Phi_1(t)$ is an integral function in $R$ which satisfies the condition
\[ P(\tau) \Phi_1(\tau) \equiv 0 \pmod{\mathfrak p \mathfrak p_1 \mathfrak p_2 \ldots} \]
By Theorem 4. it follows from this that
\weq{28}{\Phi_1(\tau) \equiv 0 \pmod{\mathfrak p_1 \mathfrak p_2 \ldots},}
%CONT>>
